\chapter{Combined oral contraceptive pill}
\index{Combined oral contraceptive pill}

\medskip
\noindent\textit{\textbf{Under review.} This chapter is an AI-assisted educational draft; it carries no source citations and is not a substitute for professional medical advice.}

\section*{Definition and Overview}
The combined oral contraceptive pill (the pill, or COCP) is a medicine taken by mouth each day to prevent pregnancy. It contains two female hormones -- an oestrogen and a progestogen -- which together stop ovulation, thicken the mucus at the opening of the womb, and make the lining of the womb unsuitable for implantation. When taken correctly, the pill is one of the most effective reversible methods of contraception. It is also used to treat period problems and other hormone-related conditions. The pill does not protect against sexually transmitted infections (STIs), so condoms are often recommended as well.

\section*{Epidemiology}
The combined pill is among the most widely used contraceptive methods in the world. In many high-income countries, roughly one in five to one in four women of reproductive age who use contraception choose the pill, and the proportion is similar across most age groups of childbearing women. Because the pill is very effective when taken reliably, missed doses are a common cause of the unintended pregnancies that do occur. The pill is also frequently prescribed for menstrual disorders, which affect a substantial minority of women. Its use declined in some countries in the late twentieth century because of concern about blood-clot risk, but it remains a first-line contraceptive for most healthy women.

\section*{Aetiology and Pathophysiology}
The pill works by preventing pregnancy rather than by treating a disease. Its effects arise from the two hormones it contains:

\begin{itemize}
    \item \textbf{Stopping ovulation:} the oestrogen and progestogen together suppress the hormones from the brain that trigger the ovary to release an egg, so most cycles produce no egg
    \item \textbf{Thickening cervical mucus:} the progestogen makes the mucus at the neck of the womb thicker and stickier, which blocks sperm from reaching the womb
    \item \textbf{Thinning the womb lining:} the hormones make the lining of the womb thin and less suitable for a fertilised egg to implant
    \item \textbf{Packet schedule:} the pill is taken for 21 days followed by 7 pill-free days (or dummy pills), during which a withdrawal bleed, similar to a light period, occurs
\end{itemize}

When pills are missed, taken late, or absorbed poorly (for example with vomiting or diarrhoea), the hormonal levels fall enough for ovulation to return, and the risk of pregnancy rises sharply.

\section*{Clinical Features}
There are no symptoms of the pill itself; rather, it produces expected changes and possible side effects:

\begin{itemize}
    \item A predictable, regular withdrawal bleed in the pill-free week, usually lighter and less painful than a natural period
    \item Lighter, shorter, and more regular periods over time
    \item Common early side effects that usually settle within a few months: breast tenderness, nausea, headaches, and spotting between periods
    \item Mood changes or reduced libido in some women
    \item An absence of bleeding is not necessarily a problem, but a missed withdrawal bleed can also indicate pregnancy
    \item The benefits are experienced as reduced period pain, less heavy bleeding, and improved acne in some women
\end{itemize}

\section*{Differential Diagnosis}
When a woman on the pill has symptoms, other causes must be considered:

\begin{itemize}
    \item \textbf{Pregnancy:} missed pills or a missed withdrawal bleed raise the possibility of pregnancy
    \item \textbf{Breakthrough bleeding:} spotting from missed pills, interactions with other medicines, or early use vs.\ a gynaecological cause such as infection or polyp
    \item \textbf{Headaches:} pill-related headache vs.\ migraine with aura, which affects the safety of the pill
    \item \textbf{Chest pain or calf swelling:} blood clot (thrombosis) must be excluded promptly in a woman on the pill
    \item \textbf{Breast symptoms:} pill-related tenderness vs.\ breast disease requiring assessment
    \item \textbf{Unrecognised conditions:} raised blood pressure or liver problems may be uncovered at routine review
\end{itemize}

\section*{When to Seek Medical Care}
See a doctor or sexual health clinic if you are considering the pill, and arrange a regular review once you are taking it. Seek advice early if:

\begin{itemize}
    \item You miss two or more pills, or start a packet late, and have had sex recently
    \item You have vomiting or severe diarrhoea lasting more than a day
    \item You develop severe headaches, especially new migraine with visual disturbance
    \item You have unexpected or persistent bleeding between periods
    \item You think you may be pregnant
\end{itemize}

\textbf{Contact your local emergency services or seek urgent care for:}
\begin{itemize}
    \item Sudden chest pain or difficulty breathing (possible blood clot in the lung)
    \item A painful, swollen, or hot calf or thigh (possible deep vein thrombosis)
    \item Sudden weakness, numbness, or difficulty speaking (possible stroke)
    \item Severe abdominal pain
    \item Sudden visual loss or double vision
\end{itemize}

\section*{Diagnosis and Investigations}
Before starting the pill, and during routine review, a clinician assesses suitability:

\begin{itemize}
    \item \textbf{History:} age, smoking, weight, blood pressure, migraines, and personal or family history of blood clots, heart disease, breast or liver disease
    \item \textbf{Blood pressure measurement:} the pill can raise blood pressure, so it is checked before and during use
    \item \textbf{Pregnancy test:} if there is any chance of pregnancy
    \item \textbf{Assessment of migraine:} distinguishing migraine with aura from migraine without aura, since aura makes the pill unsafe
    \item \textbf{STI screening:} offered because the pill does not protect against infection
    \item \textbf{Follow-up review:} at 3--6 months and then regularly, checking blood pressure, side effects, and satisfaction
\end{itemize}

\section*{Management}
The pill is prescribed and managed by a doctor or nurse, with an individualised approach:

\begin{itemize}
    \item \textbf{Starting:} begin the pill on the first day of a period (or within the first 5 days) so it works immediately; otherwise extra precautions are needed for 7 days
    \item \textbf{Taking it:} one pill at roughly the same time each day; a missed pill is taken as soon as remembered, and 2 or more missed pills may require emergency contraception and a week of condom use
    \item \textbf{Choice of formulation:} the dose of oestrogen and type of progestogen may be adjusted to reduce side effects
    \item \textbf{Continuous use:} running packets together without the pill-free week, or tricycling, can reduce periods and is safe when advised
    \item \textbf{Drug interactions:} some medicines, including certain antibiotics, epilepsy drugs, and St John's wort, reduce effectiveness and require extra precautions
    \item \textbf{Switching:} moving to a different pill, another hormonal method, or an implant or IUD if side effects persist
    \item \textbf{Stopping:} fertility returns quickly after stopping, and there is no need for a break for health reasons
\end{itemize}

\section*{Complications}
The pill is very safe for most healthy women, but it carries genuine risks that are small:

\begin{itemize}
    \item Blood clots in the leg (deep vein thrombosis) or lung (pulmonary embolism); the absolute risk is small but roughly 2--4 times that of non-users, and higher in the first year and in smokers
    \item A small increase in the risk of stroke and heart attack, mainly in smokers, older women, and those with high blood pressure or migraine with aura
    \item A small increase in the risk of breast and cervical cancer, with some protection against ovarian and endometrial (womb lining) cancer
    \item Worsening of migraine, mood changes, and reduced libido in some women
    \item Gallbladder problems in susceptible women
    \item Unintended pregnancy with a low risk of ectopic pregnancy, if pills are missed
\end{itemize}

\section*{Prognosis and Outcomes}
When taken correctly, the pill is more than 99\% effective at preventing pregnancy, and its effectiveness is around 91\% with typical use, meaning roughly 9 in 100 women a year become pregnant when pills are sometimes missed. Most side effects settle within the first few months, and many women use the pill for years without problems. Fertility returns to normal within days to weeks of stopping. The small risks of blood clots, stroke, and some cancers are balanced against the benefits of reliable contraception, and for most healthy women the pill is one of the safest and most effective options available.

\section*{Prevention and Self-Care}
\begin{itemize}
    \item Take the pill at the same time each day and keep a spare packet to avoid running out
    \item Start a new packet on time and never take more than a 7-day break
    \item Use condoms in addition to the pill to protect against STIs
    \item Use extra precautions if you miss pills, have vomiting or diarrhoea, or start interacting medicines
    \item Do not smoke while taking the pill, as smoking greatly increases the risk of blood clots and heart problems
    \item Have your blood pressure checked regularly
    \item Carry emergency contraception and condoms for missed-pill situations
    \item See a clinician promptly for unexplained bleeding, severe headaches, or leg or chest symptoms
\end{itemize}
